\subsection{The Countermodule of a Semisimple Module}

Let A be a ring. Denote the set of classes of simple A-modules by S. For every \(\lambda \in S\) , choose a simple A-module \(S_{\lambda}\) of class \(\lambda\) , and denote the opposite ring of the field of endomorphisms of \(S_{\lambda}\) by \(D_{\lambda}\) . We view \(S_{\lambda}\) as an \((A, D_{\lambda})\) -bimodule.

Let \(M\) be a semisimple A-module, and let \(B\) be the endomorphism ring of \(M\) . Denote the bicommutant of \(M\) by \(C\) . For every \(\lambda \in \mathcal{S}\) , denote the left \((D_{\lambda}, B)\) -bimodule \(\mathrm{Hom}_A(S_{\lambda}, M)\) by \(V_{\lambda}\) . Finally, denote the support of the A-module \(M\) by \(\mathcal{S}_M\) (VIII, p.~66); it is also the set of elements \(\lambda\) of \(\mathcal{S}\) such that \(V_{\lambda}\) is nonzero.

Remark 1. --- The canonical description \(\alpha_{M}\) of the A-module M is an isomorphism of left (A, B)-bimodules. By the corollary of Proposition 9, VIII, p.~71, the mapping \(f \mapsto (\operatorname{Hom}(1_{\mathrm{S}_{\lambda}}, f))_{\lambda \in \mathcal{S}_{\mathrm{M}}}\) from B to \(\prod_{\lambda \in \mathcal{S}_{\mathrm{M}}} \operatorname{End}_{\mathrm{D}_{\lambda}}(\mathrm{V}_{\lambda})\) is a ring isomorphism.

PROPOSITION 5. --- a) The countermodule of M is semisimple.

\begin{enumerate}
\def\labelenumi{\alph{enumi})}
\setcounter{enumi}{1}
\item
  For every \(\lambda \in \mathcal{S}_{\mathrm{M}}\) , the B-module \(V_{\lambda}\) is simple, and its commutant is equal to \((D_{\lambda})_{V_{\lambda}}\) .
\item
  The mapping \(\lambda \mapsto \mathrm{cl}(\mathrm{V}_{\lambda})\) is a bijection from the support of the A-module M to the support of its countermodule.
\item
  For every \(\lambda \in \mathcal{S}_{\mathrm{M}}\) , the B-submodule \(\mathrm{M}_{\lambda}\) is the isotypical component of type \(\mathrm{V}_{\lambda}\) of the B-module M, and the multiplicity of \(\mathrm{V}_{\lambda}\) in M is equal to \(\dim_{\mathrm{D}_{\lambda}}(\mathrm{S}_{\lambda})\) .
\item
  For \(s \in S\) , denote the mapping \(\varphi \mapsto \varphi(s)\) from \(V_{\lambda} = \operatorname{Hom}_{\mathrm{A}}(\mathrm{S}_{\lambda}, \mathrm{M})\) to M by \(\widetilde{s}\) . It is B-linear. The resulting mapping \(s \mapsto \widetilde{s}\) from \(S_{\lambda}\) to \(\operatorname{Hom}_{\mathrm{B}}(\mathrm{V}_{\lambda}, \mathrm{M})\) is an isomorphism of \((\mathrm{A}, \mathrm{D}_{\lambda})\) -bimodules.
\end{enumerate}

Let \(\lambda \in \mathcal{S}_{\mathrm{M}}\) . Denote the ring \(\mathrm{End}_{\mathrm{D}_{\lambda}}(\mathrm{V}_{\lambda})\) by \(\mathrm{E}_{\lambda}\) ; since \(\mathrm{V}_{\lambda}\) is a nonzero \(\mathrm{D}_{\lambda}\) -vector space, it is a simple \(\mathrm{E}_{\lambda}\) -module, (VIII, p.~45, Example 3), and its commutant is equal to \((\mathrm{D}_{\lambda})_{\mathrm{V}_{\lambda}}\) (VIII, p.~82, Corollary 1 of Theorem 2). Since \(\mathrm{E}_{\lambda}\) is the ring of homotheties of the B-module \(\mathrm{V}_{\lambda}\) (VIII, p.~71, Corollary of Proposition 9), this proves b).

The canonical description \(\alpha_{M}\) of M defines an isomorphism \(\alpha_{\lambda}\) from \(V_{\lambda}\otimes_{D_{\lambda}^{\circ}}S_{\lambda}\) to \(M_{\lambda}\) . Since \(V_{\lambda}\) is a simple B-module, the B-module \(V_{\lambda}\otimes_{D_{\lambda}^{\circ}}S_{\lambda}\) is isotypical of type \(V_{\lambda}\) (VIII, p.~61, Proposition 1); the same therefore holds for the B-module \(M_{\lambda}\) , which proves a).

By Remark 1 above, there exist elements \(e_{\lambda}\) of B, for \(\lambda\) running through \(\mathcal{S}_{\mathrm{M}}\) , such that \((e_{\lambda})_{\mathrm{V}_{\lambda}} = 1_{\mathrm{V}_{\lambda}}\) and \((e_{\lambda})_{\mathrm{V}_{\mu}} = 0\) for \(\mu \in \mathcal{S}_{\mathrm{M}}\) with \(\mu \neq \lambda\) . The simple B-modules \(V_{\lambda}\) are therefore pairwise nonisomorphic, which proves c) and the first assertion of d). The B-module \(M_{\lambda}\) is isomorphic to \(V_{\lambda} \otimes_{D_{\lambda}^{\circ}} S_{\lambda}\) , so \(\dim_{D_{\lambda}}(S_{\lambda})\) is the multiplicity of \(V_{\lambda}\) in M (II, §3, No.~7, p.~255, Corollary 1).

The mapping \(\sum_{\lambda\in\mathcal{S}_{M}}\alpha_{\lambda}\) from \(\bigoplus_{\lambda\in\mathcal{S}_{M}}V_{\lambda}\otimes_{D_{\lambda}^{\circ}}S_{\lambda}\) to M provides a description (VIII, p.~69, Definition 5) of the semisimple B-module M. By VIII, p.~70, Proposition 8, b), for every \(\lambda\in\mathcal{S}_{M}\) , the mapping from \(S_{\lambda}\) to \(\mathrm{Hom}_{\mathrm{B}}(\mathrm{V}_{\lambda},\mathrm{M})\) described in e) is bijective and \(D_{\lambda}\) -linear. Since it is obviously A-linear, this proves e).

Remark 2. --- It follows from the proof that the mapping

\[
\sum_ {\lambda \in \mathcal {S} _ {\mathrm{M}}} V _ {\lambda} \otimes_ {D _ {\lambda} ^ {\mathrm{o}}} S _ {\lambda} \longrightarrow M
\]

induced by the canonical description of M is a description of the countermodule of M.

PROPOSITION 6. --- a) Viewed as an (A, B°)-bimodule, M is semisimple.

\begin{enumerate}
\def\labelenumi{\alph{enumi})}
\setcounter{enumi}{1}
\item
  For every \(\lambda \in \mathcal{S}_{\mathrm{M}}\) , \(\mathrm{M}_{\lambda}\) is a simple (A, B°)-sub-bimodule of M.
\item
  For every \((\mathrm{A},\mathrm{B}^{\circ})\) -sub-bimodule N of M, there exists a unique subset \(\Lambda\) of \(\mathcal{S}_{\mathrm{M}}\) such that N is equal to \(\oplus_{\lambda \in \Lambda}\mathrm{M}_{\lambda}\) .
\end{enumerate}

Let \(\lambda\) be in \(\mathcal{S}_{\mathrm{M}}\) . The left A-module \(S_{\lambda}\) and the right B-module \(V_{\lambda}\) are simple, and \(D_{\lambda}\) is the opposite ring of the field of endomorphisms of \(S_{\lambda}\) . By Corollary 2 of VIII, p.~63, the (A, B°)-bimodule \(S_{\lambda} \otimes_{D_{\lambda}} V_{\lambda}\) is simple, and the same holds for \(M_{\lambda}\) , which is isomorphic to it. This proves b), and a) follows.

If \(\lambda\) and \(\mu\) are distinct in \(S_{M}\) , then \(M_{\lambda}\) and \(M_{\mu}\) are not isomorphic as A-modules, nor a fortiori as \((A,B^{\circ})\) -bimodules. Assertion c) follows by Corollary 2 of VIII, p.~68.

PROPOSITION 7. --- a) For every element c of the bicommutant C of M and every \(\lambda\in\mathcal{S}_{M}\) , there exists a unique element \(c_{\lambda}\) of \(\operatorname{End}_{\mathrm{D}_{\lambda}}(\mathrm{S}_{\lambda})\) such that for every \(\varphi\in\operatorname{Hom}_{\mathrm{A}}(\mathrm{S}_{\lambda},\mathrm{M})\) and every \(s\in S_{\lambda}\) , we have \(c\varphi(s)=\varphi(c_{\lambda}s)\) .

\begin{enumerate}
\def\labelenumi{\alph{enumi})}
\setcounter{enumi}{1}
\item
  Endow the \(S_{\lambda}\) , for \(\lambda\) running through \(S_{M}\) , with the C-module structure defined by a). Then the canonical mapping \(\alpha_{M}\) from \(\bigoplus_{\lambda\in S_{M}}S_{\lambda}\otimes_{D_{\lambda}}V_{\lambda}\) to M is an isomorphism of \((\mathrm{C},\mathrm{B}^{\circ})\) -bimodules.
\item
  The mapping \(c \mapsto (c_{\lambda})_{\lambda \in \mathcal{S}_{\mathrm{M}}}\) is an isomorphism from C to \(\prod_{\lambda \in \mathcal{S}_{\mathrm{M}}}\mathrm{End}_{\mathrm{D}_{\lambda}}(\mathrm{S}_{\lambda})\) .
\end{enumerate}

Assertions a) and c) follow from VIII, p.~71, Proposition 9 because the canonical mapping \(\alpha_{\mathrm{M}}\) from \(\bigoplus_{\lambda \in \mathcal{S}} (\mathrm{S}_{\lambda} \otimes_{\mathrm{D}_{\lambda}} \mathrm{W}_{\lambda})\) to M provides a description of the B-module M (Remark 2). Moreover, \(\alpha_{\mathrm{M}}\) is (C, B°)-linear, which proves b).

Remark 3. --- Endow \(S_{\lambda}\) , for \(\lambda \in S_{M}\) , with the C-module structure given by Proposition 7, a). If we replace A with B and B with C in Proposition 5 (VIII, p.~85), then we see that for every \(\lambda \in S_{M}\) , the left C-module \(S_{\lambda}\) is simple, with commutant \(D_{\lambda}\) , that the isotypical component of type \(S_{\lambda}\) of the C-module M is equal to \(M_{\lambda}\) , and that the mapping \(\lambda \mapsto \mathrm{cl}_{\mathrm{C}}(\mathrm{S}_{\lambda})\) is a bijection from the support of the A-module M to the support of the C-module M. Finally, observe that the A-linear mappings and C-linear mappings from \(S_{\lambda}\) to M are identical, as are the A-submodules and C-submodules of M, and that the rings \(\operatorname{End}_{\mathrm{A}}(\mathrm{M})\) and \(\operatorname{End}_{\mathrm{C}}(\mathrm{M})\) are equal.

Let M be a semisimple A-module. Denote by Z the center of the bi-commutant C of the A-module M; it is also the center of the commutant B of M. Endow M and the \(S_{\lambda}\) , for \(\lambda \in S_{M}\) , with Z-module structures deduced by restriction of scalars from the C-module structures. For every \(\lambda \in S_{M}\) , denote the center of the field \(D_{\lambda}\) by \(Z_{\lambda}\) .

PROPOSITION 8. --- a) The mapping \(z \mapsto (z_{\mathrm{S}_{\lambda}})_{\lambda \in \mathcal{S}_{\mathrm{M}}}\) is an isomorphism from Z to the product of the fields \(\mathbf{Z}_{\lambda}\) .

\begin{enumerate}
\def\labelenumi{\alph{enumi})}
\setcounter{enumi}{1}
\item
  The A-module M is isotypical and nonzero if and only if Z is a field.
\item
  Let \(\Lambda\) be a subset of \(\mathcal{S}_{\mathrm{M}}\) . Denote by \(e_{\Lambda}\) the unique element of \(\mathrm{Z}\) such that \((e_{\Lambda})_{\mathrm{S}_{\lambda}} = 1_{\mathrm{S}_{\lambda}}\) for \(\lambda \in \Lambda\) and \((e_{\Lambda})_{\mathrm{S}_{\lambda}} = 0\) for \(\lambda \in \mathcal{S}_{\mathrm{M}} - \Lambda\) . We have \((e_{\Lambda})_{\mathrm{M}_{\lambda}} = 1_{\mathrm{M}_{\lambda}}\) for \(\lambda \in \Lambda\) and \((e_{\Lambda})_{\mathrm{M}_{\lambda}} = 0\) for \(\lambda \in \mathcal{S}_{\mathrm{M}} - \Lambda\) .
\item
  If the support \(S_{M}\) of M is finite, then the mapping \(\Lambda \mapsto e_{\Lambda}Z\) is a bijection from the set of subsets of \(S_{M}\) to the set of ideals of Z, and the mapping \(a \mapsto aM\) is a bijection from the set of ideals of Z to the set of \((A,B^{o})\) -sub-bimodules of M. These bijections are isomorphisms of ordered sets. The reverse bijection sends an \((A,B^{o})\) -sub-bimodule N of M to the ideal consisting of the elements z of Z that send N into M.
\end{enumerate}

For \(\lambda \in \mathcal{S}_{\mathrm{M}}\) , \(Z_{\lambda}\) is the common center of the commutant \(D_{\lambda}\) and the bi-commutant \(C_{\lambda}\) of the A-module \(S_{\lambda}\) . By Proposition 7, c) above, the mapping \(c \mapsto (c_{\lambda})_{\lambda \in \mathcal{S}_{\mathrm{M}}}\) is an isomorphism from C to \(\prod_{\lambda \in \mathcal{S}_{\mathrm{M}}} C_{\lambda}\) . By restriction to the centers, we obtain the isomorphism \(z \mapsto (z_{S_{\lambda}})_{\lambda \in \mathcal{S}_{\mathrm{M}}}\) from Z to \(\prod_{\lambda \in \mathcal{S}_{\mathrm{M}}} Z_{\lambda}\) , whence a).

The ring \(\prod_{\lambda \in \mathcal{S}_{\mathrm{M}}}Z_{\lambda}\) is a field if and only if the set \(\mathcal{S}_{\mathrm{M}}\) has a single element, whence b).

Assertion c) follows from Proposition 7, a). Suppose that \(\mathcal{S}_{\mathrm{M}}\) is finite. It follows from a) and Proposition 8 of I, §8, No.~10, p.~109, that the mapping \(\Lambda \mapsto e_{\Lambda}\mathrm{Z}\) is an isomorphism of ordered sets from \(\mathfrak{P}(\mathcal{S}_{\mathrm{M}})\) to the set of ideals of Z. Let \(\Lambda\) be a subset of \(\mathcal{S}_{\mathrm{M}}\) . By c), we have the relation \(e_{\Lambda}\mathrm{ZM} = e_{\Lambda}\mathrm{M} = \bigoplus_{\lambda \in \Lambda}\mathrm{M}_{\lambda}\) ; because of Proposition 6, c) of VIII, p.~86, it remains to describe the inverse bijection. But \(z \in \mathrm{Z}\) sends M into \(\bigoplus_{\lambda \in \Lambda}\mathrm{M}_{\lambda}\) if and only if \(z = e_{\Lambda}z\) , that is, \(z \in e_{\Lambda}\mathrm{Z}\) .

COROLLARY. --- Suppose that A is an algebra over an algebraically closed commutative field K and that M is a semisimple A-module that is finite-dimensional as a vector space over K. For every \(\lambda\) in \(S_{M}\) , denote by \(e_{\lambda}\) the projector of M with image \(M_{\lambda}\) and kernel \(\oplus_{\lambda\neq\mu}M_{\mu}\) . Then \((e_{\lambda})_{\lambda\in\mathcal{S}_{M}}\) is a basis of the vector space Z over K.

Since M is a finite-dimensional vector space over K that is the direct sum of the family of nonzero submodules \((\mathrm{M}_{\lambda})_{\lambda\in\mathcal{S}_{\mathrm{M}}}\) , the set \(S_{M}\) is finite, and each of the spaces \(S_{\lambda}\) , for \(\lambda\in S_{M}\) , is finite-dimensional over K. Since the field K is algebraically closed, we have \(D_{\lambda}=Z_{\lambda}=K\) (VIII, p.~47, Theorem 1), and the mapping \(z\mapsto(z_{\mathrm{S}_{\lambda}})_{\lambda\in\mathcal{S}_{\mathrm{M}}}\) is an isomorphism from Z to \(K^{S_{M}}\) (Proposition 8, a)). The corollary then follows from part c) of Proposition 8.

